\mainmatter

\chapter{IL GIORNO DI COLOMBO}\label{il-giorno-di-colombo}

\hfill\break

I ruhar ci colpirono nel Giorno di Colombo. Ogni paese aveva un nome per
il giorno in cui i ruhar attaccarono; quello che fece presa, dopo un
po', fu Giorno di Colombo. Immagino abbia senso. Eccoci là, a
vagabondare ignari nel cosmo sulla nostra piccola biglia blu, come i
nativi americani nel 1492. Da oltre l'orizzonte giunsero navi di una
società aggressiva e avanzata dal punto di vista tecnologico\ldots{} e
bam! Addio bei vecchi tempi in cui noi umani ci ammazzavamo gli uni gli
altri. Ecco perché il Giorno di Colombo. Calza.

Quando il cielo del mattino scintillò, di quella che in seguito avremmo
appreso essere la nave da guerra dei ruhar che saltava in alta orbita,
eravamo incuriositi. Quando centrali elettriche, raffinerie, fabbriche e
altri siti industriali in giro per il pianeta iniziarono a essere
colpiti dall'orbita dai dardi dei cannoni a rotaia supersonici, eravamo
scioccati. Per quanto mi riguarda, quando un mezzo di trasporto da
combattimento ruhar cadde dal cielo e strisciò attraverso un campo di
patate fuori dalla mia minuscola città natale nel Nord del Maine, quello
fu il momento in cui cominciai a essere ufficialmente allarmato. Era una
mattina presto d'inizio ottobre, la festa del Giorno di Colombo in
America. Ero a casa in licenza dall'esercito, in visita dai miei. Mi
prendevo una pausa dopo che il mio battaglione era tornato a casa dal
servizio di pace in Nigeria. Che lavoro di merda. Ero felice di essere a
casa negli States. Mise fine al mio congedo una scia di fuoco che solcò
il cielo, perché il mezzo di trasporto ruhar passò dritto sopra il mio
furgone. Disegnò un arco sul lago e si schiantò nel campo di patate di
Olafsen, arando tuberi e terra, finché la sua prua finì mezza sommersa
in uno stagno. Oscillò da un lato all'altro per un minuto, con un suono
urlante proveniente dai motori, poi decollò e volò basso e instabile
sopra la linea degli alberi, verso il centro della città, rilasciando
una scia di fumo dalla parte inferiore.

I nostri presunti alleati, i kristang, pensano che il mezzo di trasporto
d'assalto ruhar sia stato ostacolato durante la caduta dall'orbita e
quindi abbia mancato di poco il bersaglio designato. Questa è l'unica
ragione possibile per cui i ruhar avrebbero invaso Thompson Corners, nel
Maine. Cazzo, non volevo più stare là, ecco il motivo principale per cui
mi ero arruolato nell'esercito. La mia città natale è abbastanza
graziosa, solo che non è nulla di speciale. Ci sono patate e bestiame e
pecore e altri allevamenti, certo, e alcune persone, come mio padre,
lavoravano alla grande cartiera giù a Milliconack. Potevi tirare avanti
producendo un po' di legname, lavorare come guida di caccia o di pesca,
fare il saldatore di tanto in tanto, qualunque cosa. Nessuno nel Nord
del Maine vive di un solo lavoro. Comunque, Thompson Corners non è il
genere di obiettivo strategicamente vitale che verrebbe in mente a dei
pianificatori militari che dovessero decidere il punto in cui far cadere
dall'orbita un mezzo di trasporto da combattimento, con una dozzina di
soldati armati fino ai denti. I soldati ruhar, carini pelosi baffuti
bastardi quali sono, senza dubbio aspettarono finché il loro mezzo di
trasporto non slittò frenando sul prato di fronte alla scuola elementare
nel centro della città, aprirono le porte, rimasero ammirati alla
magnifica vista di Thompson Corners e chiesero al pilota dove diavolo
fossero. I soldati sono soldati, che abbiano pelliccia, pelle o scaglie.
Perciò, com'era logico che fosse, i ruhar lanciarono un missile
sull'edificio più imponente della zona, il magazzino delle patate, e lo
divelsero in modo impressionante. Intendo che lo fecero saltare in aria
di brutto, quei soldati dovevano avere qualcosa contro le patate. Poi
distrussero i due ponti sul fiume Scanicutt: il ponte ferroviario e il
vecchio ponte autostradale in cemento risalente agli anni Trenta, quando
era stato costruito dalla Works Progress Administration di Franklin
Delano Roosevelt. Avevamo avuto un sacco di pioggia e il fiume era molto
gonfio, perciò, col ponte fuori uso, l'unico modo che avevo per entrare
in città era guidare fino a Woodford e attraversare là il corso d'acqua.
Com'era quel vecchio modo di dire del New England: non puoi arrivare là
restando qua, giusto? Buona idea, se solo le strade non fossero state
intasate di traffico congestionato e un centinaio di altre persone non
avessero avuto la stessa idea che avevo avuto io, nello stesso momento.

Quando vidi quella nave d'assalto arrivare, lasciandosi dietro una scia
di fumo, diretta verso Thompson Corners, ero già nel furgoncino dei miei
genitori, diretto in città, per andare a prendere mia sorella a casa
della sua amica. Dieci minuti dopo che il primo sito era stato colpito,
la radio annunciò che il governatore aveva appena dichiarato lo stato di
emergenza e aveva esortato tutti a restare calmi, poi ogni comunicazione
s'interruppe. Niente radio, niente cellulari, niente Tv, niente
elettricità. Non avevo bisogno di aspettare istruzioni, stavo andando a
prendere la mia sorellina per tornare a casa e resistere con la mia
famiglia, quando mi resi conto di cosa diavolo stava succedendo. Dietro
il sedile c'erano il fucile da caccia di mio padre, una scatola di
proiettili buoni per le quaglie e non molto altro. Il che dimostra
quanto avessi le idee chiare quella mattina. Tutta la mia attrezzatura
militare era a Fort Drum, nello stato di New York. Ero, dopotutto, in
congedo. Superando la collina, vidi le macerie del ponte e quasi mi
scontrai con una fila di auto che stava cercando di entrare in città
come me. Bill Geary, un vigile del fuoco volontario e capitano in
pensione della guardia nazionale del Maine, stava cercando di
organizzare la gente per passare oltre Woodford ed entrare a Thompson
Corners usando il vecchio sterrato. Io, come un idiota, gridai che il
mio furgoncino aveva quattro ruote motrici. Quasi tutti nel Nord del
Maine avevano quattro ruote motrici, fossero pure quelle di una vecchia
Subaru malridotta. Poiché ero l'ultimo della fila, mi fecero girare per
primo e tre tizi, la cui auto era finita in un fosso, saltarono sul mio
furgone. Ripartimmo rombando, come in un film, quando arriva la
cavalleria.

Prima che ci districassimo nel traffico in uscita per attraversare il
ponte a Woodford e prima che fossimo rimbalzati sullo sterrato malmesso,
a metà strada verso la città, i ruhar avevano già messo in sicurezza il
centro di Thompson Corners. La città era vuota perché nessun umano aveva
atteso un invito su carta bollata per svignarsela da lì. Il vicesceriffo
era stato preso dal panico e aveva sparato un paio di colpi con la sua 9
millimetri di servizio, finché i ruhar non si erano seccati e, con
quello che sembrava un razzo anticarro, avevano fatto saltare in aria la
stazione della Shell in cui si nascondeva.

Da allora ho visto i ruhar da vicino un sacco di volte. No, non mangiano
gli umani. E no, non uccidono i bambini. Credete alla propaganda, se
volete, io so cosa ho visto. Se il vicesceriffo non avesse fatto fuoco,
forse i ruhar non avrebbero ucciso una sola persona in città. Non potevo
biasimarli: se un idiota mi avesse sparato addosso, gli avrei dato fuoco
con un razzo anch'io. Lo so, perché avevo fatto la stessa cosa in
Nigeria.

In ogni caso, il cancello del servizio forestale oltre la strada
sterrata era chiuso. Perdemmo cinque minuti mentre un tizio, tre camion
davanti a noi, provava a far saltare la serratura con un fucile. Il
servizio forestale aveva previsto che qualcuno ci avrebbe provato, ma
quel lucchetto non si smosse. Un tizio speronò il cancello e lo
distrusse, assieme al radiatore del suo camion, che poi dovette essere
rimosso a spinta dalla strada, prima che noialtri potessimo infilarci e
passare oltre.

Sì, sì, tutti hanno una versione di quel giorno, questa è la mia, perciò
zitti e ascoltate. Una cosa che ho imparato è che i soldati
dell'esercito ruhar sono come i soldati del resto della galassia:
vogliono finire il combattimento e tornarsene alle loro caserme, o tane,
nel loro caso. Li odiavo? Cazzo, sì, ma non penso che avessero
intenzione di uccidere delle persone, se non come danno collaterale
almeno. Qualunque fosse il loro obiettivo sulla Terra, quei soldati
avevano mancato la loro zona di atterraggio e stavano cercando di
sfruttare al meglio la situazione. Sarebbe stato un bene per tutti se
quei ruhar se ne fossero stati seduti sul loro mezzo di trasporto guasto
a grattarsi le palle, avessero chiamato la versione ruhar del soccorso
stradale e avessero aspettato un carro attrezzi. Le operazioni di
combattimento non funzionano così. Qualcosa s'incasina in ogni missione,
ma tu ti adatti e fai del tuo meglio per raggiungere il tuo obiettivo.
Questo gruppo di ruhar decise che il suo obiettivo era mettere in
sicurezza la contea di Penobscot, che avesse senso o meno. I kristang ci
dissero che il piano più probabile dei ruhar era distruggere la nostra
infrastruttura industriale e riportarci all'età della pietra, così non
avremmo rappresentato una minaccia per loro. Se quello era il loro
obiettivo originario, lo avevano mancato di qualche centinaio di
chilometri, quando erano atterrati a Thompson Corners. I kristang
stavano per lo più dicendoci la verità rispetto al motivo per cui i
ruhar ci avevano colpiti, sebbene mentissero su tutto il resto, come
mostrerò alla fine.

Non avremmo neppure dovuto combattere i ruhar, non erano loro i nostri
nemici. Lo erano i nostri alleati.

Ma sarà meglio che cominci dal principio.

\hfill\break

Mi chiamo Joe Bishop, avevo vent'anni quando i ruhar attaccarono la
terra ed ero uno specialista dell'esercito americano. Prima che
l'esercito me li tagliasse, avevo capelli un po' più lunghi della norma,
di un indeterminato colore biondo-marrone che avevo preso da mia madre.
Lei lo chiamava ``marrone topo'' e si tingeva i capelli di biondo dorato
da quando ho memoria. Gli occhi azzurri li ho presi da entrambi i miei
genitori e il mio metro e ottantasette di altezza viene senz'altro da
parte di padre, mia madre arriva a stento a uno e cinquanta senza
scarpe. Al liceo giocavo in terza base nella squadra di baseball, ero
ricevitore distante nella squadra di football e guardia tiratrice di
riserva nella squadra di basket, anche se l'ultimo anno ho abbandonato
il basket. La verità è che non ero una star né del baseball né del
football. Lavoravo duro, mettevo la squadra al primo posto e vincevamo
la nostra razione di partite. Quando era giunto il momento di mandare le
domande d'iscrizione ai college, io non sapevo dove volessi andare, o
cosa volessi fare da grande. L'unica cosa che sapevo era che non volevo
rimanere seduto dietro un banco tutto il giorno. E che volevo andarmene
da Thompson Corners. Mio padre era stato nell'Aeronautica per un paio
d'anni, poi era entrato nelle riserve come meccanico. Stesso tipo di
lavoro che faceva alla cartiera. Gli piaceva lavorare con le mani,
aggiustare cose, e in un certo senso piaceva anche a me. I soldi
scarseggiavano e io non volevo seppellirmi di debiti con prestiti
studenteschi; perciò la carriera militare mi suonava bene.

Quando mi ero arruolato nell'esercito, l'avevo fatto perché volevo
servire il mio paese e perché il compenso mi sarebbe servito per pagare
il college. Anche quel genere di vita mi attraeva, mi piaceva stare
all'aria aperta: campeggio, caccia, pesca, trekking, canoa.
L'addestramento era stato duro, certo, niente che non mi aspettassi, ed
ero orgoglioso di aver superato le prove di base ed essere stato
assegnato dove volevo: la X divisione di fanteria di montagna a Fort
Drum, New York. La missione di pace in Nigeria non era quello che avrei
desiderato, ma ci era stato ordinato e ci ero andato. Mi aveva sorpreso
che il mantenimento della pace comportasse l'uccisione di così tante
persone, ma è così e basta.

Ora sapete, quindi, perché me ne stavo sdraiato sotto un cespuglio in
cima a una piccola cresta che sovrastava il centro della mia città
natale, fissando quella nave da trasporto ruhar guasta. Sdraiato lì a
cercare di capire cosa, se mai fosse stato necessario, avremmo dovuto
fare.

«Quello è un grosso criceto maledetto, Bish, non ti sbagli», disse Tom
Paulson, mentre mi restituiva il binocolo. «Cosa faremo?»

«Non lo so ancora. Fammici pensare.» C'erano molti veterani militari
nella mia città, ne avevo uno con me, ma Tom era stato un impiegato
della Marina vent'anni prima e io ero un fante dell'esercito con
esperienza di combattimento recente e in servizio attivo. Immagino fosse
normale per gli altri aspettarsi che fossi io ad avere delle idee.
Sapevano che ero stato in combattimento in Nigeria, ma scontrarsi con i
fanatici della disorganizzata milizia locale nella boscaglia era del
tutto diverso dall'azzuffarsi con giganti criceti spaziali nella mia
città natale nel Nord del Maine.

«Cazzo, dov'è tutto l'armamentario quando ne hai davvero bisogno? È
tutto quello che abbiamo?» Guardai con sgomento la collezione di fucili
da caccia, fucili e la stramba 9 mm. Tutti a Thompson Corners avevano
una pistola, perché tutti andavano a caccia o, se non altro, avevano
bisogno di tenere gli orsi lontani dalle loro mangiatoie per uccelli.
«Andiamo, nessuno ha un vecchio M60 in soffitta? Forse un Ak?»

«Cazzo, Bish, voglio uccidere un alce per mangiarmelo, mica vaporizzare
quel cazzo di coso», disse Tom. «Le licenze di caccia non te le tirano
dietro.»

«Va bene, va bene, hai ragione.» Guardai di nuovo nel binocolo il fumo
scuro proveniente dal magazzino delle patate che, dall'altra parte della
città, vicino ai binari della ferrovia, stava ancora bruciando. E la
navicella, o mezzo da sbarco o nave d'assalto della fanteria (o comunque
la chiamassero i criceti), appariva tozza e brutta e potente, con la
prua fracassata e un'ala piegata e una sottile striscia di fumo bianco
proveniente dalla pancia, piazzata sul prato davanti alla scuola
elementare.

Criceti. Abbiamo altri nomi per loro: ratti, donnole, roditori, ma con
quella pelliccia fine e dorata, le facce rotonde e i baffi, ciò a cui
più assomigliano sono i criceti. Solo che i criceti non sono alti un
metro e ottanta, non stanno in piedi su due gambe, non indossano
giubbotti antiproiettile, caschi e occhiali, non trasportano fucili
dall'aspetto minaccioso e non scendono dall'orbita su una nave
d'assalto. Almeno credo\ldots{} Intendo dire, non ho mai avuto un
criceto, quindi cosa volete che ne sappia? Non sapevamo che fossero
criceti, finché uno di loro all'ingresso della nave − forse il pilota,
perché non aveva un fucile − non si tolse il casco, prese qualcosa dalla
tasca e iniziò a mangiarlo. Gli altri gli urlarono contro e gli rimisero
il casco, ma non prima che noi avessimo visto la sua testa pelosa e le
sue orecchie da criceto. Non erano proprio criceti, ma qualcosa di
abbastanza simile.

Susie Tobin alzò l'occhio dal mirino del suo fucile da caccia. «E la
cava? Hanno la dinamite, giusto?» Susie era alta un metro e un tappo,
come primo lavoro faceva l'insegnante alla scuola media regionale, e
guardandola non avresti mai detto che fosse in grado di sollevare il suo
vecchio fucile residuato dell'esercito. Ma avevo visto le rastrelliere
piene di corna di cervo, attaccate sul lato sud del suo fienile, quindi
di sicuro sapeva come usarlo.

«Che ce ne facciamo della dinamite?», chiese beffardo Diego. «Corriamo
verso la loro nave e la lanciamo attraverso la porta? Non faresti in
tempo a fare cento metri prima che ti abbattano.»

«Susie ha ragione», dissi, pensando all'inferno che ci avevano fatto
passare con gli ordigni esplosivi improvvisati in Nigeria. Erano un
problema sulle strade, ma pericolosi soprattutto quando pattugliavamo un
villaggio, dove le linee di avvistamento erano limitate e c'erano molti
posti per nascondere una bomba. Pattugliare, come stavano facendo i
criceti. Coppie di criceti stavano controllando edifici nel centro della
città, che, trattandosi di quello di Thompson Corners, non è grande. Si
potrebbe pensare che, essendo la festa del Giorno di Colombo, ci fossero
molti turisti in città quando i ruhar fecero un salto per una visita.
Questo per chi non ha mai avuto la fortuna di visitare la mia città
natale. Nel New England il Giorno di Colombo è una grande festa dedicata
all'osservazione delle foglie, un fine settimana in cui la gente di
città, dal profondo Sud, va in campagna con l'auto per vedere il
fogliame colorato sugli alberi, soggiornare in piccole locande
pittoresche e bed and breakfast e scattare alle foglie un sacco di foto.
Noi accogliamo turisti nella nostra parte dei Great North Woods, ma non
per sbirciare le foglie. Ora, nel Giorno di Colombo, molti degli alberi
qui intorno hanno già perso le foglie. Inoltre, la nostra parte del
Maine è comunque piuttosto piatta e, in aggiunta, i pini non cambiano
colore. La gente viene qui per il canottaggio, la pesca, la caccia e la
motoslitta. Non è ancora piena stagione per tutto questo la prima metà
di ottobre, per cui la città era piuttosto vuota quando i ruhar si
schiantarono a terra, essendo giorno di vacanza da scuola e tutto il
resto. Per fortuna, direi. Non riuscivo a pensare a cosa sarebbe
successo se la scuola elementare fosse stata piena quando i ruhar
presero il controllo della città. Non dovevo immaginare quello scenario,
l'avevo già visto nel film Red Dawn. Intendo la versione bella di quando
era in carica Reagan, però, non quella schifosa che hanno fatto dopo.

Comunque, ebbi un'idea. «Susie, tuo padre lavora alla cava, giusto? Tu e
Tom procurateci dinamite, detonatori, cavi, qualsiasi cosa ci serva per
far esplodere qualcosa a distanza.» Non avevo idea di come fare, non
avevo mai visto un candelotto di dinamite. «Diego, tu rimani qui e
controlla i criceti, vedi dove pattugliano, specie quando passano dal
vecchio municipio», divenuto la sede dell'unico ristorante della città e
di un'agenzia assicurativa. «Stan, Deb, vedete se laggiù riuscite a
trovare un camion o un furgone che possiamo far partire, abbastanza
grande da farci stare un paio di persone sul retro, qualcosa di coperto,
non un letto aperto. Restate da questa parte della città, non provate ad
attraversare la Route 11, o i criceti vi vedranno. Se trovate qualcosa,
portatela in Red Brook Road e lasciatela lì. Ci rivediamo qui.»

«Ok, Bish», annuì Tom. «Cosa faremo?»

Diedi un'altra occhiata ai ruhar. Notai la disposizione delle loro
truppe a casaccio, esaminando dove erano di pattuglia in giro per la
città e dove avevano posizioni difensive intorno all'astronave dissi:
«Ce la prenderemo comoda».

\hfill\break

«Non ci posso credere», sparai.

«Bish», disse Stan sulla difensiva, «questo è il meglio che siamo
riusciti a trovare: o questo, o un vecchio Dodge Neon.»

«La maggior parte della gente ha preso l'auto e se n'è andata in fretta,
sembra, non sono rimasti molti mezzi da questa parte del fiume»,
aggiunse Debbie.

«Sì, ma\ldots»

Tom scosse la testa: «Non esiste che viaggiamo su questo coso. Stiamo
combattendo un'invasione aliena. Non possiamo andarci su questo affare».

Dovetti concordare. Tranne su un punto: era un furgone perfetto. Avevano
trovato un furgone per le consegne, tipo quelli di FedEx, ed era in
ottime condizioni meccaniche; le gomme non erano lisce, c'erano dei
punti arrugginiti attorno ai passaruota, ma niente di serio, aveva
grandi portelli posteriori, in modo che molte persone potessero entrare
o uscire in fretta, e dicevano che funzionasse e si muovesse bene, senza
cinghie cigolanti o freni stridenti. L'avevano trovato dietro il garage
dei fratelli Davis, dove dovevano essere in corso i lavori di
riparazione dell'interno, perché nell'abitacolo tutti i rivestimenti
erano stati tolti, a parte quello del sedile del conducente. Era un buon
furgone, tranne che per quell'unico particolare.

Barney.

Barney, il grande dinosauro viola del cartone animato, con quel suo
eterno sorriso idiota. Barney, Puffi, Topolino, unicorni e un sacco di
altri personaggi immaginari erano dipinti sul furgone. Chi aveva deciso
quali personaggi ritrarre aveva fatto scelte interessanti. Tipo: perché
Iron Man stava facendo un cenno di saluto ai Puffi? Ed era Darth Vader
quello laggiù, vicino al parafango anteriore destro, oppure qualcuno
aveva iniziato a stendere la mestica nera per coprire una macchia di
ruggine e aveva deciso soltanto dopo di dar fondo alla sua vena
creativa? La maggior parte dei personaggi erano stati dipinti male, e mi
ci volle un minuto per capire che quello che pensavo fosse un Buddha
seduto doveva essere Winnie the Pooh. Winnie il Buddha? Era un furgone
dei gelati, comunque, nel caso non l'aveste ancora indovinato. Un
ridicolo furgone dei gelati, che ero abbastanza sicuro non avesse il
permesso di utilizzare nessuno dei caratteri registrati dipinti sul
lato.

Inoltre, al posto del noto marchio ``Mister Softee'', questo aveva il
marchio ``Super Softie'', mal modellato. Come faceva il gelato a essere
super-morbido? Era sciolto?

È chiaro che si trattava di un furgone dei gelati pirata, del tipo che
immaginavo vagare furtivo nei quartieri periferici di una grande città,
vendere gelati scaduti e cercare di evitare le autorità locali. Era
impossibile che il Barney viola gigante che copriva entrambi i lati e il
Barney di peluche legato al cofano passassero inosservati. Chiunque
fosse il proprietario di quel furgone doveva davvero, davvero amare
Barney.

«Potremmo spalmarci una rapida mano di vernice», suggerì Stan.

«Non abbiamo tempo da perdere a ridipingere quel cazzo di coso»,
ringhiai.

Neanche a me piaceva l'idea, ma quel furgone era tutto quello che
avevamo per andare in battaglia. «E poi i criceti non sanno chi è
Barney, potrebbero pensare che sia un feroce predatore.» A questa non
credevo neppure io.

«Oh, per l'amor di Dio!», Susie era esasperata. «Voi uomini idioti avete
paura di andare in giro in un furgone con su scritto ``morbido''? Datevi
una calmata. Guiderò io quel cazzo di coso. Joe, qual è il piano?»

\hfill\break

«Via, via, via!», gridai insieme a Stan, mentre saltavamo per stenderci
sul pianale del furgone dei gelati. Susie non se lo fece ripetere due
volte, diede gas mentre eravamo ancora a mezz'aria. Stan sarebbe
scivolato fuori dal retro se Deb non l'avesse preso per il bavero della
camicia. Tom calciò dentro i piedi di Stan e chiuse i portelli sul
retro, proprio mentre il furgone rimbalzava su una grande buca. Tom
atterrò di culo sopra il soldato ruhar. Il criceto grugnì facendoci
capire che era ancora vivo, cosa di cui non ero del tutto sicuro quando
lo avevamo catturato.

Ripensandoci, era un piano stupido ed eravamo stati fortunati. Avevo
notato che i criceti stavano attraversando il vicolo tra il ristorante e
il negozio di ferramenta lì vicino. Mentre noi eravamo fuori a fare
provviste, Diego aveva osservato e confermato che le pattuglie di
criceti passavano attraverso quel vicolo con regolarità, osservando
dalle finestre e bussando alle porte aperte per sbirciare all'interno.
Lo stavano facendo in tutta la città, e io avevo scelto la tavola calda
perché c'era una breve strada sterrata che portava dal fiume fin sul
retro del locale. La strada attraversava i boschi e aveva una buona
copertura per la maggior parte del percorso. Avevamo avuto un sacco di
pioggia da nord-est la settimana prima, perciò il fiume era quasi allo
stadio da inondazione primaverile. Era davvero ruggente mentre passava
sopra le rocce e sotto il ponte, faceva abbastanza rumore da coprire
quello del furgone in retromarcia dietro la tavola calda. Avevamo
piazzato della dinamite all'interno del locale, contro il muro esterno,
e l'avevamo fatta esplodere quando i soldati criceti erano a metà del
vicolo.

Nessuno di noi sapeva quanta dinamite usare, quindi ne avevamo usata
troppa, e aveva fatto crollare la maggior parte del muro nel vicolo.
Anche il furgone dei gelati era stato colpito da mattoni vaganti. Non
importava, la nostra bomba improvvisata aveva fatto il suo dovere: aveva
fatto uscire di testa quei due criceti e li aveva sepolti sotto i
mattoni. Tom, Stan e io avevamo catturato il criceto più vicino a noi,
che era anche quello meno sommerso dai mattoni. Debbie ci aveva coperto
le spalle con un fucile, mentre lo tiravamo fuori dal vicolo: Tom e Stan
l'avevano afferrato per un braccio e io l'avevo tenuto per i piedi,
avevamo arrancato lungo il vicolo, inciampando sui mattoni, e avevamo
gettato il criceto nel retro del furgone dei gelati.

Era un piano stupido. Diego osservava la situazione dalla collina con un
walkie-talkie che Susie aveva preso dalla cava, aveva l'altra unità nel
furgone. Diego aveva riferito che, non appena la dinamite era esplosa,
una mezza dozzina di criceti si era precipitata fuori dalla loro
astronave. Altri trenta secondi e ci avrebbero catturati. E se per
qualche motivo, visto che nessuno di noi sapeva davvero come usare la
dinamite, la nostra bomba non fosse esplosa, quei due criceti avrebbero
raggiunto a piedi il retro del vicolo e avrebbero visto il furgone che
prima non c'era. Non mi entusiasmava il calcolo delle nostre probabilità
di riuscita, contro soldati alieni dotati di giubbotti antiproiettile e
armi avanzate.

La fortuna era con noi quel giorno, anche se Tom cadde di nuovo, questa
volta contro il pulsante che controllava l'impianto stereo del furgone
dei gelati, e Turkey in the Straw risuonò dagli altoparlanti sul
tettuccio mentre il furgone rimbalzava e slittava sulla strada sterrata
lungo il fiume. «Spegni quella cazzo di roba!», urlò Susie a Stan, aveva
già il suo bel da fare a guidare. La cintura di sicurezza per il
conducente non c'era: era uno degli oggetti che avevano rimosso quando
avevano tolto i rivestimenti all'interno del furgone, e i piedi di Susie
raggiungevano a stento il pedale del gas. In più, aveva quel babbeo sul
pavimento e ci stavamo muovendo.

Stan pigiò sui pulsanti e la musica cambiò in Camptown Races, poi in Pop
Goes the Weasel, poi in un paio di canzoni a tema videogiochi di cui non
saprei dire il nome, prima che Stan stoppasse finalmente quel delirio.
Non mi sono mai sentito così assolutamente idiota in vita mia cercando
di tenere fermo un soldato alieno nel retro di un furgone dei gelati,
sbattendo col mento sul pianale di metallo, mentre musica demenziale
esplodeva dagli altoparlanti e un sorriso da maniaco percorreva la
faccia del gemello malvagio di Barney da un orecchio all'altro.

Il nostro furgone sfrecciò attraverso un tratto di strada non alberata
lungo il fiume, uno spazio libero di circa cinquanta metri. È là che
pensai che i criceti ci avrebbero sparato: avevano la visuale libera per
almeno un paio di secondi. Il motivo per cui non lo fecero, credo, è che
a quel punto i criceti che avevano raggiunto il vicolo bombardato si
erano accorti della scomparsa di uno di loro e avevano capito che era
nel nostro furgone. In ogni caso, ce la cavammo: il nostro mezzo
attraversò il tratto di strada non alberata, poi una lieve salita tra
noi e il centro della città, quindi Susie si mise in piedi sui freni per
rallentare prima di una curva. Dopo di che, rimbalzammo su una strada
asfaltata e Susie diede gas a manetta. Una volta che fummo sul fondo
asfaltato, con una guida più fluida, io e Tom legammo le mani del
criceto dietro la schiena, gli assicurammo le gambe l'una con l'altra e
io iniziai a togliergli l'attrezzatura. Lo mettemmo sotto quattro
giubbotti di piombo, di quelli che usano i dentisti quando ti fanno i
raggi X, fu una mia idea, per bloccare il segnale di qualsiasi
dispositivo di geolocalizzazione i soldati alieni potessero avere. Il
giubbotto antiproiettile attorno al busto aveva un meccanismo di sgancio
rapido, come mi aspettavo, perché un soldato, che fosse umano o alieno,
doveva potersi mettere e togliere l'attrezzatura in fretta. Gli aprii
anche il gancio del sottogola del casco e, per la prima volta, potei
vedere bene in faccia il nemico. Aveva solo un occhio aperto, gli usciva
sangue da un taglio sulla guancia, ma non mi parve fosse ferito in modo
grave, per lo più stordito e disorientato. Indossava un auricolare e un
microfono che strappai, e sulla cintura c'era una specie di radio. Dissi
a Tom di buttare tutto fuori dal finestrino, saremmo tornati a prenderlo
più tardi se ne avessimo avuto l'opportunità. La nostra priorità in quel
momento era catturare un soldato alieno in modo che i nostri militari
potessero studiarlo, vedere che tipo di nemico stavamo affrontando.

Secondo il piano, Susie ci trasportò per un chilometro e mezzo fino a
casa di Tom: lui aveva un fienile e lei guidò il furgone dritto
attraverso la porta aperta, scivolando fino a fermarsi. Per un istante
ci sentimmo tutti della serie ``Sogno o son desto?''. Nessuno di noi
poteva crederci. Il criceto si muoveva, perciò Tom ci si sedette sopra e
io gli puntai una pistola in faccia. Questo lo calmò all'istante. È
probabile che non avesse mai visto una Sig Sauer prima, ma riconobbe
un'arma da fuoco quando se la trovò a un palmo dal naso.

«Non riesco più a sentire Diego», riferì Deb tenendo il suo
walkie-talkie. «L'ultima cosa che ha detto è che gli alieni stavano
correndo qua e là, ma la loro nave non si muoveva.» L'ultima istruzione
che avevo dato a Diego era stata di lasciare la città, nella direzione
opposta, non appena avesse perso di vista il nostro furgone. Ero certo
che fosse al sicuro, sapeva come muoversi nei boschi.

«Stanno aspettando un carro attrezzi», ipotizzai.

«Cosa?», chiese Susie, come se non mi avesse sentito bene.

«Un carro attrezzi», spiegai. «La loro nave è rotta, devono aver
chiamato la loro flotta di sopra, e stanno aspettando che qualcuno venga
a riparare il guasto, o a prenderli. Portiamo Faccia Pelosa qui nella
cantina, prima che i suoi amici si organizzino e vengano a cercarlo.»

La casa di Tom aveva un vecchio seminterrato e, quando dico vecchio,
voglio dire che la sua casa risale al 1848 e il seminterrato esisteva
già prima di allora. Il seminterrato era stato ampliato e trasformato in
un bunker dalla famiglia che possedeva la fattoria negli anni Cinquanta,
ora Tom e sua moglie Margie lo usavano come deposito. Portammo il
criceto giù per le scale e liberammo spazio al centro del rifugio per
lui, Tom chiuse con un lucchetto una catena attorno alla sua caviglia e
ne fissò l'altra estremità attorno a un tubo che fuoriusciva dal
pavimento. Con il criceto comodo al centro, non c'era molto spazio per
noialtri a causa di tutti gli scaffali presenti. Vi ho detto che la
moglie di Tom, Margie, amava la frutta e le verdure? La donna aveva un
piccolo problema con l'inscatolamento: il posto era tutto pieno di
barattoli. La maggior parte delle persone in questa parte del bosco sa
fare conserve, è un modo per avere buon cibo durante l'inverno ed è
molto più economico che comprare cibo in negozio. I miei producevano
mele e marmellate, pomodori in scatola, fagioli e tutto quello che
potevamo ricavare dal nostro giardino. Il bunker di Tom faceva pensare
che Margie stesse progettando di invitare a cena l'intera X divisione e
volesse far avanzare un sacco di cibo.

«E adesso?», chiese Stan, guardando con avidità un barattolo di
marmellata di more.

«Adesso aspettiamo. I militari devono per forza arrivare presto.»
Eravamo tutti sorpresi di non vedere neanche un caccia o un elicottero
in cielo, il che mi diceva che il nemico aveva la supremazia aerea
totale. «Se tutto va bene, gli amici di questo tizio», indicai il nostro
prigioniero, «partiranno presto. Quando la nostra cavalleria arriverà,
lo consegneremo.»

``Cosa ti fa pensare che stiano per andarsene?'' chiese Deb.

``Perché gli invasori alieni non avrebbero mai messo Thompson Corners
sulla loro lista di destinazione iniziale. Questi tizi sono qui solo
perché la loro nave è rotta. La mia grande preoccupazione è che restino
nei dintorni a cercare Faccia Pelosa dopo l'arrivo del carro attrezzi.
Io resto con lui, voi andate a sud e vedete se riuscite a trovare la
guardia nazionale o la polizia di stato.''

Susie era indignata. «Stai mandando le donne a mettersi in salvo?» Aveva
le mani sui fianchi, sapevo cosa significava.

«Tu hai una famiglia.»

«Anche Tom», fece notare Susie.

«Tom ha ricevuto un messaggio da Margie», dissi sulla difensiva, «lei e
i bambini stanno bene, diretti a casa di sua madre. E la moglie di Stan
è a Portland.»

«E mio marito è giù a Milliconack e i miei figli sono a Bangor con mia
sorella. Io resto qui», disse Susie in tono empatico e sollevò il fucile
con fare enfatico.

«Bene», ammisi senza lottare molto, in realtà non volevo restare da
solo. «Tu e Stan avete dei fucili, salite nei boschi dietro la casa dove
potete coprire la strada. Deb, sali sulla roccia dietro la casa, la
conosci, vero? Dì a Susie e Stan se vedi arrivare qualcosa da questa
parte. Io e Tom staremo qui con Faccia Pelosa.»

\hfill\break

Faccia Pelosa non era loquace. Tom e io cercammo di comunicare a gesti
ma, indicandoci a dito, il criceto se ne stava seduto con la faccia di
pietra. Solo quando scivolò in una posizione scomoda, allora fece una
smorfia. Aveva sangue, sangue rosso, sul fianco. Se non altro il sangue
alieno era rosso e non verde o blu o qualcosa di davvero strano. Mi
avvicinai piano, a mani aperte per dimostrare che non avevo un coltello.
«Dobbiamo vedere dove sei ferito», dissi a ritmo lento e a voce alta,
che è sempre il modo migliore per farsi capire dagli stranieri.
Ricordate, gli stranieri sono stupidi, quindi dovete parlare a voce alta
e a ritmo lento in modo che vi capiscano. Il criceto si allontanò da me
il più possibile, si appoggiò contro uno scaffale sul muro esterno.
Indicai con un gesto il mio fianco, poi il suo, poi misi un dito nel
sangue sul pavimento e lo sollevai in modo che potesse vedere il liquido
rossastro. Scuotendo la testa e agitando il dito insanguinato, dissi a
ritmo lento al criceto: «Sei ferito. Questo è male».

«Bish», disse Tom in tono incerto, «sei sicuro di voler toccare il
sangue di quel coso?»

«Dev'essere più sicuro che toccare sangue umano. Hai delle bende
quaggiù?»

«Sì, il kit di pronto soccorso è nell'armadietto. Abbiamo anche dei
tovaglioli di carta, Margie li compra al supermercato.»

Il criceto non era ferito in modo grave, pulii la lesione, la lavai con
acqua sterile, il che lo fece sobbalzare di nuovo, poi spalmai della
lozione antisettica sul taglio e applicai una benda. Mi guardò in modo
diverso, dopo, come se fosse sorpreso che gli umani fossero civilizzati.
Disegnò una sorta di bacio con le labbra.

«Non è che vuole dell'acqua?», ipotizzò Tom. Margie aveva anche comprato
bottiglie d'acqua all'ingrosso; ne tirammo fuori tre, Tom e io bevemmo
dalle nostre, poi ne aprii un'altra e la portai alle labbra del criceto.
Vuotò mezza bottiglia senza staccarsi e fissandomi negli occhi in un
modo che mi spaventò. Poi abbassò lo sguardo in modo eloquente sulla sua
tasca anteriore sinistra, gesticolando con il naso. Con molta cautela,
gli aprii il lembo della tasca, che era una specie di cosa magnetica,
credo, e tirai fuori quella che sembrava una barretta energetica. Non
aveva una confezione brillante come le barre Hooah! che conoscevo grazie
all'esercito, solo un involucro di plastica verde, con una scritta
bianca in caratteri alieni. Aprii l'estremità dell'involucro di
plastica, di fronte a me. Che cosa stupida, e se fosse stato un qualche
tipo di esplosivo e fosse bastato tirare l'involucro per innescare la
granata? Non successe nulla.

La annusai con cura: odorava di zucchero e noce moscata mescolati alla
segatura. Era una barretta energetica di sicuro. Ne staccai un pezzo e
lo diedi in pasto al criceto che lo masticò per un bel po'\ldots{} Credo
che le loro barrette energetiche siano dure come le nostre. In una
decina di minuti, ne mangiò metà e bevve altra acqua, poi scosse la
testa quando gliene offrii un altro pezzo. Mi succedeva la stessa cosa
con le barrette energetiche umane: di solito ne bastava metà a stancarmi
la mascella. In un ulteriore tentativo di cooperazione interspecifica,
tenni la barretta di fronte al criceto mentre vi avvolgevo sopra
l'involucro e nascosi di nuovo la metà avanzata nella sua tasca, accanto
ad altre due barrette energetiche intere. Sorrise, o almeno ci provò. Il
sorriso, senza comprendere il contesto, avrebbe potuto essere visto come
l'atto di scoprire i denti con fare minaccioso. Il criceto aveva denti
non molto diversi da quelli umani: smalto bianco, sembrava, ma i due
anteriori superiori erano più grandi del normale, ma non di comica
grandezza come quelli di un castoro.

Ora che ero più vicino, vidi che la pelliccia sul suo volto era sottile,
gli copriva tutta la pelle, ma era più fine e più corta di una barba
umana.

«Lo senti?», chiese Tom. «Sembra che qualcuno stia gridando.» Per poter
sentire Susie, nel caso in cui avesse visto qualcosa, avevamo accostato
leggermente la porta del bunker, che in realtà tenevamo per lo più
chiusa, in modo che tutti i segnali radio di un eventuale trasmettitore
impiantato nel nostro criceto sarebbero stati bloccati. Almeno così
speravo. Tom uscì, lo sentivo gridare con Stan. Infilò la testa nella
porta. «Dicono che sia arrivato il carro attrezzi.»

Lasciammo il nostro criceto prigioniero nel bunker-scantinato-magazzino
e ci spingemmo nei boschi attraverso il cortile di Tom. «È venuto da
nord-ovest, ho visto la scia di condensazione per prima», riferì Deb. «È
volato proprio nel centro della città. Penso tu abbia ragione, Joe, i
criceti qui hanno chiesto un passaggio.»

«Spero che, qualunque sia la loro missione, sia più importante che
trovare un soldato disperso.» Se la nuova nave aliena aveva caricato i
criceti bloccati e lasciato l'area, il mio piano era caricare di nuovo
Faccia Pelosa nel furgone dei gelati e sfrecciare a sud sulla strada,
finché non avessimo trovato un'unità militare, con molta probabilità la
guardia nazionale, che avrebbe potuto prendere in custodia il nostro
prigioniero. «Saliamo sulla sporgenza rocciosa dove si riesce a vedere
qualcosa.»

Quello che vedemmo non prometteva niente di buono: la nuova nave aliena,
dello stesso tipo della prima, stava girando attorno alla città a circa
cinquecento metri. Sembrava che stesse volando sulla base di uno schema
di ricerca e i pod delle armi sulla fiancata si erano aperti, esponendo
file di missili. Ci furono un lampo di luce, un rombo di tuono e una
colonna di fumo dal centro della città. «Hanno fatto esplodere la loro
astronave distrutta», annunciai, come se non l'avessero intuito tutti.
«Interessante.»

«Perché?», chiese Deb.

«Credo significhi che non hanno intenzione di restare qui, almeno, non
qui», indicai la città. «E non vogliono che ficchiamo il naso nei loro
ritrovati tecnologici in loro assenza. Se avessero pianificato di
restare qui, avrebbero portato una squadra di riparazione.»

La nave salì, restò in volo a punto fisso, poi si diresse a bassa
velocità verso di noi. Non verso di noi, a dire il vero, ma a est, e
quindi volò lungo il fiume. «Merda.» Stan sputò a terra. «Mi sa che
ricevono un segnale da quella roba che abbiamo buttato fuori dal
furgone.»

Cercai di rassicurare la gente: «Niente panico. Ecco perché abbiamo
buttato quella roba fuori dal furgone, ed è per questo che il nostro
criceto adesso è sotto cemento e terra, dietro una porta d'acciaio». Da
quando il nostro prigioniero era nello scantinato, Stan aveva gettato le
coperte di piombo sopra il cofano e il tubo di scappamento del furgone
per ridurre la traccia termica. «A meno che non riescano a rintracciarci
nel fienile di Tom, ci metteranno un bel po' a cercare ovunque, e non
possono sapere che siamo già a quindici chilometri di strada da lì.»
Stan e Deb insistettero perché proseguissimo, invece di fermarci a casa
di Tom, per allontanarci il più possibile da Thompson Corners prima che
arrivasse una nave di salvataggio aliena. Bocciai quell'idea, perché non
sapevamo quando sarebbe arrivata un'altra nave aliena e perché non mi
fidavo del tutto del fatto che le coperte potessero bloccare un segnale
di localizzazione. Non sapevamo nemmeno se gli alieni usassero la radio,
potevano disporre di una tecnologia più avanzata. Se erano in grado di
rintracciare un segnale, che fossimo a un chilometro di distanza oppure
a ottanta, non avrebbe fatto differenza, ci avrebbero trovati. Lo
scantinato-rifugio antiatomico-qualunque cosa di Tom era perfetto per
nascondere un alieno tecnologicamente avanzato. Se il mio piano fosse
valido oppure no dipendeva da quello che avrebbero fatto i criceti.

La nave aliena seguì a regime ridotto il fiume, poi accelerò e costeggiò
l'autostrada verso sud, verso di noi. Dopo aver fatto due giri, calò
fuori visuale sotto la linea degli alberi. Per un momento, un solo un
momento di pace, le cose sembrarono quasi come in un normale giorno di
ottobre nelle campagne del Maine. La casa di Tom e Margie era a posto,
c'erano un paio di cataste di legna da ardere ben impilate in una fila
sul retro. Margie aveva legato dei gambi di mais attorno al lampione in
giardino, pronti in anticipo per Halloween, come sempre. Nel cortile
posteriore, teneva la vasca da bagno Madonna decorata con un\ldots{}

Oh, sarà meglio che spieghi cos'è una vasca da bagno Madonna ai nostri
soggetti poco acculturati. Prendete una vecchia vasca da bagno,
interrate per metà uno dei lati corti, dipingete l'esterno di bianco e
l'interno di azzurro chiaro e poi mettete una statua della Vergine Maria
sotto l'arco, come in una grotta. È una specie di santuario fatto in
casa. Al giorno d'oggi si possono comprare quelli già fatti in cemento,
ma questo è barare e Dio lo sa se sei stato pigro. La maggior parte
della gente ci mette attorno pietre o fiori; il santuario di Margie era
circondato da crisantemi dai colori vivaci. O almeno mi parvero
crisantemi, non sono un esperto di fiori. La Madonna c'era già quando la
coppia aveva comprato la casa, e Margie l'aveva ornata con i fiori.
Inoltre, la decorava a seconda delle stagioni. A Natale, Mary aveva un
cappellino da Babbo Natale. Proprio ora, per Halloween, Mary indossava
un costume da cavaliere Jedi, completo di spada laser luminosa, che Tom
aveva agganciato. Spero che Dio abbia un gran senso dell'umorismo.

Dall'altra parte della strada c'erano dei pini, un sacco di pini, e poi
campi pianeggianti e uno scorcio del fiume. Era una bella giornata di
sole. A parte le colonne di fumo provenienti dalla città, dal magazzino
di patate in rovine e dalla nave aliena bruciata. E a parte l'altra nave
aliena, che in quel momento risaliva verso il cielo, sfrecciava fino a
trecento metri e si fermava in volo a punto fisso. «Sì», dissi, «hanno
trovato quella roba aliena che abbiamo buttato fuori dal furgone e hanno
fatto scendere dei tizi a controllare. Quell'astronave lassù sta
fornendo supporto aereo ravvicinato ai tizi a terra. Tom, ci sono»,
contai sulle dita, «sei, no, sette, edifici tra qui e il punto in cui si
trovano loro?»

«Dieci», mi corresse Tom, «non hai contato il McDonald e il Burgess,
arretrato rispetto alla strada. E la vecchia bancarella di frutta della
stazione di servizio, che è abbandonata da vent'anni. Sono dieci
edifici, con strutture multiple come fienili e garage. Ispezioneranno a
fondo ciascuno di essi, giusto, e ci vorrà del tempo.»

Mi morsi il labbro, cosa che faccio quando sto pensando. «Il comandante
alieno lassù non ha idea di dove abbiamo portato il suo soldato.» Che
cosa avrei fatto io se fossi stato nei suoi panni? «Potremmo essere in
qualsiasi casa o fienile da queste parti, potremmo essere in una grotta
nel bosco, potremmo essere a una decina di chilometri lungo la strada.
Non ne ha idea.» A meno che il nostro prigioniero non avesse una specie
di trasmettitore rilevabile perfino nel bunker. Possedevano, dopotutto,
la tecnologia per viaggiare tra le stelle. Mi diedi un colpetto sul
mento, altra cosa che faccio quando sto pensando. «Hanno inviato una
sola nave. Se questa fosse una missione di ricerca e salvataggio ne
avrebbero inviate diverse. Quel tizio è sceso dall'orbita prima che
prendessimo Faccia Pelosa e aveva in mente solo di caricare i criceti
della nave distrutta.»

«Cosa significa?», chiese Susie.

«Significa che i prossimi cinque minuti ci diranno quali sono i loro
ordini. Se quella nave recupera i tizi a terra e vola via, vuol dire che
si stanno attenendo al loro piano di missione originario e si
occuperanno di ricerca e salvataggio più tardi. Se quella nave resta in
giro, vuol dire che hanno cancellato il piano originario e stanno
aspettando rinforzi per perlustrare l'area.»

«È questo che faresti?», domandò ancora Susie.

«È quello che faremmo noi, l'esercito.»

La nave aliena calò di nuovo di quota, poi salì e riprese a volare a
punto fisso più vicino a noi. «Cazzo! È l'opzione B. Ha appena fatto
scendere altri tizi per iniziare a perquisire gli edifici.» Vidi Susie
controllare di nuovo il suo fucile. «Susie, dimentica le stronzate da
cowboy di Alamo. Abbiamo solo tre fucili, un fucile da caccia e due
pistole. Se si avvicinano, ci dirigiamo a est attraverso i boschi. I
missili su quella nave potrebbero ucciderci tutti da diecimila metri di
distanza.»

«Ci arrenderemo?»

«Non ci arrendiamo, siamo realisti.» Non potevo credere di stare
discutendo di questo con Susie. «Non possiamo usare di nuovo il furgone
dei gelati, è troppo appariscente e quella nave aliena è troppo vicina.»

«Bene!», rispose Susie. «Avvolgiamo il nostro Faccia Pelosa in coperte
di piombo e portiamolo nel bosco. È soltanto a un chilometro e mezzo da
qui, in direzione del fiume, possiamo\ldots»

«Ehi!», gridò Deb. «In cielo! Altre luci!» Alzammo gli occhi e vedemmo
il firmamento mattutino brillare di nuovo. C'erano un sacco di luci
lassù. «Oh, accidenti. Stanno portando altre navi?», chiese Deb. «Non
può essere, basta!»

Distogliemmo lo sguardo, delle macchie ci nuotavano negli occhi a
seguito di un'esplosione molto luminosa nel cielo proprio sopra di noi.
Ammiccai fissando la mia ombra per terra che tremolava, mentre altre
esplosioni rischiaravano l'aria.

«Sì!» Stan scosse il pugno. «Stiamo contrattaccando! Distruggiamo i
bastardi a suon di testate nucleari!»

«Assolutamente no.» Scossi la testa. «Le nostre armi non hanno testate
con sensori a infrarossi: possono colpire solo un bersaglio
prestabilito, a terra.» A meno che l'Aeronautica non avesse dei
giocattoli incredibili di cui non ero a conoscenza. Avevo ancora la
vista offuscata, mi schermai gli occhi e guardai verso l'orizzonte.
Altre esplosioni nel cielo, più in lontananza. «Non siamo noi.» Era un
impulso elettromagnetico: gli alieni avevano fatto esplodere un missile
in cielo per distruggere i nostri dispositivi elettronici? Allora non
potevamo saperlo, ma la seconda serie di luci nel cielo era quella di un
gruppo di battaglia kristang che saltava in orbita, poi attaccava i
ruhar.

«Allora chi è?», chiese Tom.

«Ehi, guarda!» Debbie indicò la strada. La nave aliena che era sulle
tracce del nostro prigioniero stava scendendo di nuovo, in fretta. Restò
a terra solo pochi minuti per poi decollare ancora, e questa volta cabrò
e andò dritta, veloce. Sentimmo un boato sonico e la nave si lasciò
dietro una scia di condensazione. Ovunque stesse andando, aveva fretta.
«Qualsiasi cosa stia succedendo lassù, quella nave ha appena ricevuto
nuovi ordini. Rimettiamo Faccia Pelosa nel furgone e portiamolo
all'armeria della guardia nazionale prima che tornino.»

«E allora?», chiese Tom. Mentre ci arrampicavamo sul sentiero, tutti noi
guardavamo il cielo. Dovevamo allontanare Faccia Pelosa dalla zona prima
che gli alieni cambiassero di nuovo idea.

«Allora, vedo quali sono i miei ordini.» Se non fossi riuscito a
contattare la X divisione, per il momento mi sarei messo in contatto con
una guardia nazionale locale o un'unità di riserva. «E sopravviviamo. Ho
la sensazione che il magazzino di Margie sarà più importante di queste
armi, con l'inverno alle porte.»